\section{Setting and standing facts}\label{sec:setting}

All spaces are real.  We write $s(t):=\sqrt{1+t^2}$ throughout.

\subsection{The base, the blocks and Mart\'in's lemmas}
We use the canonical base on $X=c_0$, as in the accompanying note
\cite{PrA}: $H$ is a Hilbert space, $U:H\to c_0$ is compact with dense
range, and
\[
 B_q=B_{c_0}+U(B_H),\qquad q^*(a)=\|a\|_1+\|U^*a\|_H\quad(a\in\ell_1=X^*).
\]
Since $U$ has dense range, $U^*$ is injective; it is compact, so
$\|U^*e_j^*\|_H\to0$.  The bidual ball is $B_{q^{**}}=B_{\ell_\infty}+U(B_H)$
(the sum of a weak$^*$-compact set and a norm compact set is weak$^*$
compact, contains $B_q$, and is contained in its weak$^*$ closure).  The
norm $q$ is a smoothing of the supremum norm in the sense of
Debs, Godefroy and Saint-Raymond \cite{DGS}.

Mart\'in's construction \cite{Martin} rests on two lemmas, which we recall
in the form stated there.

\begin{remark}[Mart\'in's Lemma A]\label{rem:lemmaA}
Let $\Phi$ be a sequence of positive numbers with $\|\Phi\|_1<1$ and let
$D_\Phi:\ell_2\to\ell_1$ be the diagonal operator with entries $\Phi(n)$.
The norm $|\cdot|_\Phi$ on $\ell_1$ with unit ball
$B_{\ell_1}+D_\Phi(B_{\ell_2})$ satisfies:
(a) $|f|_\Phi^*=\|f\|_\infty+\|D_\Phi f\|_2$ for $f\in\ell_\infty$;
(b) $|f|^*_\Phi\geq\|f\|_\infty\geq(1-\|\Phi\|_1)|f|^*_\Phi$;
(c) $|x|_\Phi\leq\|x\|_1\leq(1+\|\Phi\|_1)|x|_\Phi$;
(d) the dual norm is strictly convex, so $|\cdot|_\Phi$ is smooth;
(e) if $|f|^*_\Phi=1$, $\ip{f}{x}=|x|_\Phi$ and $|x(n)|>\Phi(n)|x|_\Phi$,
then $f(n)=\sgn(x(n))\|f\|_\infty$.
\end{remark}

\begin{remark}[Mart\'in's Lemma B and the choice of $T$]\label{rem:lemmaB}
Lemma B of \cite{Martin} reads: \emph{Let $Y$ be a separable operator
range.  Then there is a norm-one injective operator
$T:\ell_1(\N\times\N)\to Y$ such that for every $m$ the set
$\{T(e_{n,m})/\|T(e_{n,m})\|:n\in\N\}$ is dense in $S_Y$.}
Its proof takes a bijection $\sigma:\N\times\N\to\N$, a sequence
$(w'_n)$ in $S_Y$ such that every point of $S_Y$ is an accumulation point
of $(w'_n)$, and repeats the proof of \cite[Proposition~2.8]{KLMW20} with
the sequence $w_{\sigma(n,m)}=w'_n$.  Two features matter below.  First,
the sequence $(w'_n)$ is \emph{arbitrary} subject to the accumulation
property, so ``Mart\'in's space'' is a family of spaces indexed by the
admissible choices of $T$.  Second, every block $m$ follows the same
sequence $(w'_n)$, so that near duplicates $Te_{n,m}/\|Te_{n,m}\|\approx
Te_{n,m'}/\|Te_{n,m'}\|$ across blocks are built into the construction (up
to the perturbations produced in \cite{KLMW20}).  No other property of $T$
is used in \cite{Martin}, and no other property will be used here unless
it is stated explicitly.
\end{remark}

\begin{definition}[Admissible operators]\label{def:admissible}
An operator $T:\ell_1(\N\times\N)\to X^*=\ell_1$ is \emph{admissible} if
\begin{enumerate}
\item[(T-a)] $q^*(Te_{n,m})\le1$ for all $(n,m)$, with equality for some
$(n,m)$;
\item[(T-b)] $T$ is injective;
\item[(T-c)] $Y:=T(\ell_1(\N\times\N))$ satisfies $Y\cap c_{00}=\{0\}$;
\item[(T-d)] for every $m$, the set $\{u_{n,m}:n\in\N\}$, where
$u_{n,m}:=Te_{n,m}/q^*(Te_{n,m})$, is norm dense in $S_{q^*}$.
\end{enumerate}
\end{definition}

Since $\NA((X,q),\R)=c_{00}$ (Proposition~\ref{prop:smooth}(c) below),
(T-c) is the condition $Y\cap\NA(q)=\{0\}$ of \cite{Martin}.  As $S_{q^*}$
has no isolated points, (T-d) implies that every tail
$\{u_{n,m}:n\geq n_0\}$ is dense in $S_{q^*}$.

\subsection{The norm $p$}
Fix an admissible $T$.  For $m\in\N$ put
\begin{gather*}
 \Phi_m(k):=2^{-m-k}q^*(Te_{k,m})\le 2^{-m-k},\qquad
 \lambda_{k,m}:=m\Phi_m(k),\\
 (R_mx)(k):=\lambda_{k,m}u_{k,m}(x)\qquad(x\in X).
\end{gather*}
Let $D_m$ be the diagonal operator with entries $\Phi_m(k)$, let
$|\cdot|_m$ be the norm of Remark~\ref{rem:lemmaA} for $\Phi=\Phi_m$ on
$V_m=\ell_1$, and $N_m(w):=|w|_m^*=\|w\|_\infty+\|D_mw\|_2$.  For a nonempty
set $I\subseteq\N$ of \emph{blocks} put
\begin{gather*}
 V:=\Big(\bigoplus_{m\in I}V_m\Big)_{\ell_1},\qquad
 V^*=\Big(\bigoplus_{m\in I}(\ell_\infty,N_m)\Big)_{\ell_\infty},\\
 Lx:=(R_mx)_{m\in I},\qquad p(x):=q(x)+\|Lx\|_V .
\end{gather*}
For $I=\N$ this is Mart\'in's norm (with the canonical base); for
$I=\{1,\dots,N\}$ we obtain the finite-block norms $p_N$ of \cite{PrB}.
Since $|(R_mx)(k)|\le m2^{-m-k}q(x)$, $\|R_m\|\leq m2^{-m}$; each $R_m$ is a
norm limit of finite-rank operators, hence compact, and $L$ is compact.
For $w\in V^*$,
\begin{equation}\label{eq:Lstar}
 L^*w=\sum_mR_m^*w_m=T\Big(\sum_{m\in I}\sum_k m2^{-m-k}w_m(k)\,e_{k,m}\Big)\in Y,
\end{equation}
an absolutely convergent series; in particular $L^*$ is injective on
$V^*$.  For $\theta\in X^{**}=\ell_\infty$,
$(R_m^{**}\theta)(k)=\lambda_{k,m}\theta(u_{k,m})$, so
$R_m^{**}(X^{**})\subseteq\ell_1$.

\begin{remark}[Finite and infinite block sets]\label{rem:blocks}
By the remark on Mart\'in tails in \cite{PrB}, if the set
$\NA((c_0,p_N),F)$ is dense in $\mathcal L((c_0,p_N),F)$ for arbitrarily
large $N$, then
$\NA((c_0,p),F)$ is dense for Mart\'in's $p$ ($I=\N$): one has
$p=p_N+s_N$ with $s_N\le\eta_N\,q\le\eta_Np_N$, $\eta_N\to0$, and the
lift lemma of \cite{PrB} applies.  Hence finite block sets suffice for the
density problem.  Sections~\ref{sec:setting}--\ref{sec:averaging} are
valid for every nonempty $I$ (certificates involve finitely many blocks);
the one place where $I=\N$ needs an extra hypothesis is
Corollary~\ref{cor:decomposition}.  Sections~\ref{sec:second}
and~\ref{sec:resonance} assume that $I$ is finite.
\end{remark}

\subsection{Elementary structure}

\begin{lemma}[Dual ball]\label{lem:dualball}
$B_{p^*}=B_{q^*}+L^*(B_{V^*})$, and for every $h\in X^*$
\[
 p^*(h)=\min\{\max(q^*(A),\|W\|_{V^*}):\ A+L^*W=h\},
\]
the minimum being attained.  Moreover $p^{**}(\theta)=q^{**}(\theta)+
\|L^{**}\theta\|_V$ for $\theta\in X^{**}$, and $L^{**}(X^{**})\subseteq V$.
\end{lemma}

\begin{proof}
If $q^*(A)\le1$ and $\|W\|\le1$ then $(A+L^*W)(x)\le q(x)+\|Lx\|=p(x)$.
The set $B_{q^*}+L^*(B_{V^*})$ is convex and weak$^*$ compact ($L^*$ is
weak$^*$-continuous), and its support function on $X$ is
$q(x)+\|Lx\|_V=p(x)$, which is also the support function of $B_{p^*}$; two
weak$^*$-closed convex sets with the same support function on $X$
coincide.  The formula for $p^*$ follows by homogeneity and compactness.
Taking suprema over $B_{p^*}$ gives the bidual formula; $L^{**}$ maps into
$V$ because $L$ is compact.
\end{proof}

\begin{lemma}[Block threshold {\cite[Lemma~5.6]{PrB}}]\label{lem:threshold}
Fix a block and drop the index $m$.  Let $\zeta\in\ell_1\setminus\{0\}$ and
$w=J(\zeta)$ its norming functional ($N(w)=1$, $w(\zeta)=|\zeta|$).  Put
$M=\|w\|_\infty$, $C=\|Dw\|_2$ and $P=\{k:|w(k)|=M\}$.  Then $M>0$, $C>0$,
$M+C=1$, and
\[
 \zeta/|\zeta|=\alpha+D^2w/C,\qquad \|\alpha\|_1=1,\quad
 \supp\alpha\subseteq P,\quad \sgn\alpha(k)=\sgn w(k)\ (k\in\supp\alpha).
\]
In particular $P\ne\emptyset$; off $P$ one has
$w(k)=C\zeta(k)/(\Phi(k)^2|\zeta|)$, and $|\zeta(k)|>|\zeta|\Phi(k)^2M/C$
implies $k\in P$ and $w(k)=M\sgn\zeta(k)$.
\end{lemma}

\begin{proof}
$w\neq0$ and $D$ is injective, so $M,C>0$.  The ball
$B_{\ell_1}+D(B_{\ell_2})$ is closed ($D$ is compact), so
$\zeta/|\zeta|=\alpha+D\beta$ with $\|\alpha\|_1,\|\beta\|_2\le1$.  Then
$1=\ip{w}{\alpha}+\ip{Dw}{\beta}\le M\|\alpha\|_1+C\|\beta\|_2\le M+C=1$,
so $\ip{w}{\alpha}=M=M\|\alpha\|_1$ and $\beta=Dw/C$; the first equality
forces $\|\alpha\|_1=1$ and $\alpha$ to live on $P$ with the signs of $w$.
The remaining claims follow by reading the identity coordinatewise.
\end{proof}

For a block vector $\zeta=R_m^{**}\theta$ the condition of the last
sentence reads $|\theta(u_{k,m})|>\vartheta_m\Phi_m(k)$ with the
\emph{threshold constant} $\vartheta_m:=|\zeta|_mM_m/(mC_m)$.  For $k\in P$
we define the \emph{margin}
$\mu_{k,m}:=|\theta(u_{k,m})|-\vartheta_m\Phi_m(k)\ge0$; reading
Lemma~\ref{lem:threshold} at $k\in P$ gives the identity
\begin{equation}\label{eq:margin}
 |\zeta|_m|\alpha_m(k)|=\lambda_{k,m}\mu_{k,m}\qquad(k\in P_m).
\end{equation}
Peaks with $\alpha_m(k)=0$ (equivalently $\mu_{k,m}=0$) are called
\emph{degenerate}.  We write $Q_m:=\N\setminus P_m$ (strict non-peaks) and
$\gap_m(k):=M_m-|w_m(k)|$.

\begin{lemma}[Rigidity and independence]\label{lem:rigidity}
\textup{(a)} If $b_++L^*W_+=b_-+L^*W_-$ with $b_+-b_-\in c_{00}$ and
$W_\pm\in V^*$, then $b_+=b_-$ and $W_+=W_-$.
\textup{(b)} Let $J\subseteq\N$ be cofinite, $I$ finite, and for each
$m\in I$ let $\Lambda_m\subseteq\N$ be finite and $P_m'\subseteq\N$ be
infinite, $s\in\{\pm1\}^{P'_m}$.  Put
$\pi_m:=\sum_{k\in P'_m}s_k\lambda_{k,m}u_{k,m}$.  Then the restrictions to
$c_{00}(J)$ of the functionals $u_{k,m}$ \textup{(}$k\in\Lambda_m$,
$m\in I$\textup{)} and $\pi_m$ \textup{(}$m\in I$\textup{)} are linearly
independent; consequently $\nu\mapsto(u_{k,m}(\nu),\pi_m(\nu))$ maps
$c_{00}(J)$ onto the corresponding finite-dimensional space.
\end{lemma}

\begin{proof}
(a) $b_+-b_-=L^*(W_--W_+)\in Y\cap c_{00}=\{0\}$, and $L^*$ is injective.
(b) A vanishing combination
$\sum c_{k,m}u_{k,m}+\sum_md_m\pi_m$ on $c_{00}(J)$ is an element of $Y$
supported in the finite set $\N\setminus J$, hence zero by (T-c).  By
(T-b) its coefficient at every $e_{k,m}$ vanishes: at the infinitely many
$k\in P'_m\setminus\Lambda_m$ this coefficient is $d_ms_km2^{-m-k}$, so
$d_m=0$, and then $c_{k,m}=0$.
\end{proof}

\begin{proposition}[Strict convexity and forced decomposition]\label{prop:forced}
\textup{(a)} Each $R_m^{**}$ is injective, and $p^{**}$ is strictly convex.

\textup{(b)} Every $f\in S_{p^*}$ has a unique normer $\xi\in S_{p^{**}}$.
Put $q_0:=q^{**}(\xi)$.  Then $0<q_0<1$, $\|L^{**}\xi\|_V=1-q_0$, every
$R_m^{**}\xi\ne0$, and there is exactly one decomposition $f=a+L^*w$ with
$q^*(a)\le1$ and $\|w\|_{V^*}\le1$; it satisfies
\[
 q^*(a)=1,\quad a(\xi)=q_0,\quad w_m=J_m(R_m^{**}\xi),\quad N_m(w_m)=1
 \quad(m\in I).
\]
\textup{(c)} Put $\nu:=\|U^*a\|_H>0$ and $e:=U^*a/\nu$.  Then
$\zh:=\xi/q_0=z+Ue$ with $z\in B_{\ell_\infty}$ and $z_j=\sgn a_j$ on
$F:=\supp a$; and $a(\zh)=1$, $b(\zh)=b(z)+\ip{U^*b}{e}$ for
$b\in\ell_1$.
\end{proposition}

\begin{proof}
(a) If $\theta(u_{k,m})=0$ for all $k$, then $\theta$ vanishes on a dense
subset of $S_{q^*}$, so $\theta=0$.  Let
$p^{**}(\xi)=p^{**}(\eta)=p^{**}((\xi+\eta)/2)=1$.  Since $q^{**}$ and each
$|R_m^{**}\cdot|_m$ are seminorms whose sum is $p^{**}$
(Lemma~\ref{lem:dualball}), we get
$|R_m^{**}(\xi+\eta)|_m=|R_m^{**}\xi|_m+|R_m^{**}\eta|_m$ for every $m$.
Fix $m$; $\zeta_1=R_m^{**}\xi$ and $\zeta_2=R_m^{**}\eta$ are nonzero, and a
norming functional $w$ of $\zeta_1+\zeta_2$ norms both.  If $\xi,\eta$ were
not positively proportional there would be $v\in S_{q^*}$ with
$\xi(v)>0>\eta(v)$ (use surjectivity of $v\mapsto(\xi(v),\eta(v))$ if they
are independent).  Choose $n_j\to\infty$ with $u_{n_j,m}\to v$.  Since
$\Phi_m(n_j)\to0$, Lemma~\ref{lem:threshold} applied to $\zeta_1$ and to
$\zeta_2$ gives $w(n_j)=\|w\|_\infty$ and $w(n_j)=-\|w\|_\infty$ for large
$j$, which is absurd.  Hence $\eta=\lambda\xi$ with $\lambda>0$, and
$\lambda=1$.

(b) The normers of $f$ form a convex subset of $S_{p^{**}}$, hence a point.
For any decomposition $f=a+L^*w$ with $q^*(a),\|w\|\le1$,
$1=f(\xi)=a(\xi)+w(L^{**}\xi)\le q^{**}(\xi)+\|L^{**}\xi\|=1$, so
$a(\xi)=q_0$ and $w_m(R_m^{**}\xi)=|R_m^{**}\xi|_m$ for each $m$.  As
$R_m^{**}\xi\ne0$ and $|\cdot|_m$ is smooth, $w_m=J_m(R_m^{**}\xi)$ is
determined, hence so is $a=f-L^*w$.  Finally $q_0>0$ because $q^{**}$ is a norm and
$\xi\ne0$, $q_0<1$ because $L^{**}\xi\ne0$, and
$q_0=a(\xi)\le q^*(a)q_0$ gives $q^*(a)=1$.

(c) $\xi/q_0\in B_{q^{**}}=B_{\ell_\infty}+U(B_H)$, so $\xi/q_0=z+Uh$ with
$\|z\|_\infty,\|h\|\le1$; then $1=a(z)+\ip{U^*a}{h}\le\|a\|_1+\|U^*a\|=1$
forces $a(z)=\|a\|_1$ and $h=e$.
\end{proof}

Throughout, the \emph{forced data} of $f\in S_{p^*}$ are
$(\xi,q_0,a,w,F,z,\zh,e,\nu)$ and, for each block,
$\zeta_m:=R_m^{**}\xi$, $\sigma_m:=|\zeta_m|_m$, $M_m$, $C_m$, $P_m$,
$\alpha_m$, $Q_m$, $\gap_m$, $\vartheta_m$, $\mu_{k,m}$.  Note
$q_0+\sum_m\sigma_m=1$.  The set $K:=\{j\notin F:|z_j|=1\}$ is the set of
\emph{contacts} (kinks) and $J:=\N\setminus(F\cup K)$ the set of
\emph{free coordinates}; the \emph{room} of $j\in J$ is $1-|z_j|>0$.

